The specification fixes the intensity, and the intensity is what the model knows about the
session at time $t$.
This subsection reads two things off it:
the direction in which the flow is currently leaning,
and how that direction is assembled from the events that produced it.

\begin{defi}\label{def.intensityContrast}
 Let
 \begin{equation*}
  \upIntensity(t) := \sum_{e \, : \, \pressure_e = +1} \intensity(t),
  \qquad
  \downIntensity(t) := \sum_{e \, : \, \pressure_e = -1} \intensity(t) .
 \end{equation*}
 The \emph{intensity contrast} is
 \begin{equation}\label{eq.intensityContrast}
  \intensityContrast(t) := \upIntensity(t) - \downIntensity(t)
  = \pressure^{\transpose}\intensity[ ](t)
  = \pressure^{\transpose}\baseIntensity
  + \pressure^{\transpose}\excitation \decayedCounts(t) .
 \end{equation}
\end{defi}

The restriction of $\pressure$ to $\lbrace -1,0,+1\rbrace$ is what the middle equality needs:
for a general real vector $\pressure^{\transpose}\intensity[ ]$ is a weighted contrast and not
the difference of two sums, and $\upIntensity$, $\downIntensity$ have no meaning.

Under direction symmetry the first term of \eqref{eq.intensityContrast} vanishes, by
\eqref{eq.baselineSymmetry}, and
\begin{equation}\label{eq.contrastFromState}
 \intensityContrast(t) = \pressure^{\transpose}\excitation \decayedCounts(t) .
\end{equation}
The contrast is then a function of the state alone.

\begin{prop}\label{prop.signedExcitation}
 Let the specification be direction-symmetric, and set
 \begin{equation}\label{eq.signedExcitation}
  \signedExcitation_{\eone}
  := \frac{\big( \pressure^{\transpose}\excitation \big)_{\eone}}{\pressure_{\eone}} ,
  \qquad \eone = 1, \dots, \numEventTypes .
 \end{equation}
 Then $\signedExcitation$ takes the same value on the two members of each pair, and
 \begin{equation}\label{eq.contrastFromSignedExcitation}
  \intensityContrast(t)
  = \sum_{\eone} \pressure_{\eone} \signedExcitation_{\eone}
  \decayedCounts_{\eone}(t) .
 \end{equation}
\end{prop}
\begin{proof}
 Write $u := \excitation^{\transpose}\pressure$.
 Then $\swapMatrix u
 = \big( \swapMatrix \excitation^{\transpose} \swapMatrix \big) \swapMatrix \pressure
 = -u$ by \eqref{eq.directionSymmetry},
 so $u_{\swapMatrix(\eone)} = -u_{\eone}$;
 and $\pressure_{\swapMatrix(\eone)} = -\pressure_{\eone}$,
 so the ratio \eqref{eq.signedExcitation} is unchanged by the swap.
 Equation \eqref{eq.contrastFromSignedExcitation} is
 \eqref{eq.contrastFromState} written entry by entry.
\end{proof}

We read $\signedExcitation_{\eone}$ as the signed excitation strength of the type $\eone$:
one event of type $\eone$ raises the intensity of the types that push its own way, relative
to those that push the other way, by $\signedExcitation_{\eone}$.
A type with $\signedExcitation_{\eone} > 0$ begets its own pressure,
one with $\signedExcitation_{\eone} < 0$ begets the opposite,
and one with $\signedExcitation_{\eone} = 0$ begets neither.
The sign is a property of the kernel and not of the type:
the same six types carry either sign under different specifications,
and Section \ref{sec.forecastingTheMid} is where that matters.

\subsubsection*{The contrast and the imbalance}

The intensity contrast and the order flow imbalance are statistics of the same past events,
weighted differently.

By \eqref{eq.contrastFromState} the contrast is a function of the state $\decayedCounts$ and
of the excitation matrix.
It is available only to an observer who knows $\excitation$ and $\decay$,
and under direction symmetry it is blind to $\baseIntensity$,
so it says nothing about the immigrants that are about to arrive.
It reads the types of the past events and weights them by how long ago they occurred.

The imbalance is model free.
By Proposition \ref{prop.signedSizeContribution} it reads the same past events,
but weights them by their sizes and by nothing else:
every event in the window counts its own $\orderSize_n$, whatever its age,
and every event outside the window counts zero.
The kernel governs when events happen and of which type;
the marks govern how far they move the price.
No function of $\decayedCounts$ recovers a size.
